%Schriftart
\usepackage[utf8]{inputenc}
\usepackage[T1]{fontenc}
\usepackage[ngerman]{babel}
%Times New Roman (CZG-Vorgabe: Times New Roman oder Arial, 12 pt)
\usepackage{newtxtext}
\usepackage{romannum} %römische Ziffern in Code: \Romannum{1} (groß) \romannum{1} klein
\usepackage{tabularx}
\usepackage{booktabs}
%schriftgröße der Überschriften
\setkomafont{chapter}{\Large}
\setkomafont{section}{\large}
\setkomafont{subsection}{\normalsize}

%zeilenabstand der itemizes
\usepackage{enumitem}
\setlist[itemize]{itemsep=0pt, parsep=0pt, topsep=0pt, partopsep=0pt}

%beginn fußnoten mit leerzeichen
\deffootnote{1em}{1em}{\makebox[1em][r]{\thefootnotemark\ }}

%normale Fußnoten mit Punkt schließen lassen, wenn noch kein '.', '!', '?' vorhanden
\let\orgfootnote\footnote
\newcommand\myautodot{%
  \ifthenelse{\the\spacefactor>\sfcode`,}{}{.}%
}
\renewcommand\footnote[2][\empty]{%
  \ifx#1\empty%
    \orgfootnote{#2\myautodot}%
  \else%
    \orgfootnote[#1]{#2\myautodot}%
  \fi%
}

%Absätze werden über parskip geregelt. Die Werte sind parskip=true/false/half
%(muss vor geometry stehen, sonst werden die Seitenränder überschrieben)
\KOMAoptions{parskip=half}
%Abbildungs- und Tabellenverzeichnis im Inhaltsverzeichnis aufführen
\KOMAoptions{listof=totoc}
%keine Punkte hinter Gliederungs-/Tabellennummern (auch im Anhang: "Tabelle A.1:")
\KOMAoptions{numbers=noenddot}
%Inhaltsverzeichnis nur bis zur nummerierten Ebene 3 (x.x.x)
\setcounter{tocdepth}{2}

%anpassung seitenränder (CZG: oben, links, rechts 2,5 cm; unten 2,0 cm)
\usepackage[
    left=25mm,   % linker Rand
    right=25mm,  % rechter Rand
    top=25mm,    % oberer Rand
    bottom=20mm, % unterer Rand (bis zum Text)
    footskip=8mm % Seitenzahl ca. 1,2 cm über der Blattkante (wie Fußzeile in Word)
]{geometry}

%Festlegen des Seitenlayouts
\usepackage[onehalfspacing]{setspace} %Zeilenabstand
%alternativ \linespread{1.5}
%Kein BCOR/DIV mehr: KOMA würde damit den Satzspiegel neu berechnen und die
%geometry-Ränder überschreiben.

%Seitenzahlen in der Fußzeile zentriert, keine Kopfzeile
\usepackage[automark]{scrlayer-scrpage}
\clearpairofpagestyles
\cfoot*{\pagemark}
\pagestyle{scrheadings}

%Fußnoten fortlaufend über die ganze Arbeit nummerieren (nicht pro Kapitel)
\counterwithout{footnote}{chapter}

%Name der Verfasserin / des Verfassers in kleinerer Schrift unter der Überschrift
\newcommand{\autor}[1]{\par\noindent{\small\textit{#1}}\par}

%Abstand vor Kapitelbeginn
\RedeclareSectionCommand[
  beforeskip=0pt,
  afterindent=false
]{chapter}

%Schriftart der Überschriften serif in bold
\addtokomafont{disposition}{\rmfamily}
%Schriftart der Überschriften in serif ohne bold
%\setkomafont{disposition}{\normalfont}

%Für das Einfügen eines Glossars: siehe https://en.wikibooks.org/wiki/LaTeX/Glossary

%stützt die Schnittstelle von float und listings zu tocbasic (genutzt von KOMA-script)
\usepackage{scrhack}

%Zum Auskommentieren von größeren Blöcken
\usepackage{comment}
%Als Platzhalter
\usepackage{lipsum}
%Einfachere Verwendung von Anführungszeichen
\newcommand{\qt}[1]{\glqq#1\grqq{}}
%Angabe von Bildquellen
\newcommand{\quelle}[1]{\\\vspace{.1cm}{Quelle: #1}}
%Angabe von Tabellenquellen (unter der Tabelle, Tabellenüberschrift steht oben)
\newcommand{\tabquelle}[1]{\par\vspace{.1cm}{Quelle: #1}}

%Bilder
\usepackage{graphicx}
\usepackage{float}
\graphicspath{{./abbildungen/}}
\usepackage[export]{adjustbox}
\usepackage{array}


%Einfügen von Unformatiertem typewriter Text
\usepackage{verbatim}
%über mehrere Zeilen mit \begin{Verbatim}[breaklines=true]\end{Verbatim}
\usepackage{fancyvrb}


%Alles was das Formelherz begehrt
\usepackage{amstext} 
\usepackage{amsmath}
%newtxmath (Times-Mathematik) enthält die amssymb-Symbole und muss nach amsmath geladen werden
\usepackage{newtxmath}

%Für die Literaturverzeichnisse
\usepackage{csquotes}
%\usepackage[backend=biber, bibstyle=authoryear, citestyle=apa]{biblatex}
%Für Hauptname der Quelle in anderer Schriftart nehme man 'iso-authoryear'. (nutzt dreprecated Feld)




%Literaturverzeichnis exakt nach ISO 690-3
%Eine Alternative Erstellung von Literaturverzeichnissen über natbib (letzte Version von 2010)
%WICHTIG: bei einem Wechsel müssen alle Temporären Dateien von LaTex gelöscht werden!
%\usepackage{natbib}
%\usepackage{multibib} %Unterteilung in Internet- und Printmedien
%\newcites{online}{Internetquellenverzeichnis} 
%Statt \cite{} muss für Internetquellen nun \citeonline{} verwendet werden!



%Anklickbare Links und \url{}
\usepackage{hyperref}
%Zeilenumbrüche in URLs, standard url Paket wird von hyperref geladen, Befehl \url{} bleibt gleich
\usepackage{microtype}[final]
\urlstyle{same}

%glossar
\usepackage[acronym,hyperfirst=true,toc=true]{glossaries}
%\hypersetup{pageanchor=true}
\makeglossaries

%https://de.overleaf.com/learn/latex/Glossaries infos über Erstellung glossareintraäge und aufrufe dieser
%WICHTIG! wenn Änderungen an den Glossareinträgen vorgenommen werden,im Terminal hintereinander das hier iengeben:
%erstmal .glo datei lsöchen, dann:
%pdflatex Seminarfacharbeit.tex
%makeglossaries Seminarfacharbeit
%pdflatex Seminarfacharbeit.tex
%nur dann funktioniert das Glossar auch in der PDF-Datei
\begin{comment}
    zum kopieren:
    \newglossaryentry{}
    {}
\end{comment}

\newglossaryentry{varianzanalyse}{
  name={Varianzanalyse},
  description={Statistisches Verfahren, das prüft, ob sich die Mittelwerte mehrerer Bedingungen systematisch unterscheiden. Dazu werden die Unterschiede zwischen den Bedingungen mit den zufälligen Schwankungen innerhalb der Bedingungen verglichen.}
}

\newglossaryentry{messwiederholung}{
  name={Messwiederholung},
  description={Versuchsaufbau, bei dem dieselben Personen alle Bedingungen durchlaufen. Dadurch können Unterschiede zwischen den Personen herausgerechnet und Effekte der Bedingungen zuverlässiger erkannt werden.}
}
\newglossaryentry{trigger}{
  name={Trigger},
  plural={Trigger},
  description={Zahlencode, der während der Messung in die EEG-Aufzeichnung
  geschrieben wird, um festzuhalten, wann ein Ereignis, z.\,B. das Erscheinen
  eines Stimulus, stattfindet}
}
\newglossaryentry{claude-code}{
  name={Claude Code},
  description={KI-Assistent des Unternehmens Anthropic, der Programmcode
  schreiben und ausführen sowie Daten direkt untersuchen kann}
}
\newglossaryentry{silber-silberchlorid-elektrode}
    {
    name=Silber-Silberchlorid-Elektrode,
    first={\textit{Silber-Silberchlorid-Elektroden}},
    description={Die Silber-Silberchlorid-Elektrode ist eine häufig verwendete Elektrode in der Elektroenzephalographie und besteht aus einem Silberdraht, der mit einer Schicht aus Silberchlorid überzogen ist}
    }

\newglossaryentry{berger-effekt}
{
    name=Berger-Effekt,
    first={\textit{Berger-Effekt}},
    description={Der Berger-Effekt bezeichnet eine Blockade der Alpha-Wellen im EEG und ist nachweisbar, in dem ein Proband, der wach und entspannt ist, die Augen öffnet. Die normalerweise bei geschlossenen Augen sichtbaren Alpha-Wellen verschwinden und werden durch Beta-Wellen ersetzt. Beim Schließen der Augen kehrt der Alpha-Rhythmus zurück}
}

\newglossaryentry{brain-computer-interface}
{
    name=Brain-Computer-Interface,
    first={\textit{Brain-Computer-Interface}},
    description={Brain-Computer-Interfaces, auf deutsch Gehirn-Computer-Schnittstellen sind Systeme, die es ermöglichen, Informationen direkt aus dem Gehirn zu extrahieren und in Befehle für externe Geräte umzuwandeln}
}

\newglossaryentry{elektrolyt}
{
    name=Elektrolyt,
    first={\textit{Elektrolyt}},
    description={Ein Elektrolyt ist eine Substanz, die in der Lage ist, elektrische Ladungen zu leiten, wenn sie als Lösung oder im geschmolzenem Zustand vorliegt}
}

\newglossaryentry{analog-digital-converter}
    {name=Analog-Digital-Converter,
    first={\textit{Analog-Digital-Converter}},
    description={Ein Analog-Digital-Converter, auf deutsch Analog-Digital-Wandler, ist ein elektronisches Bauteil, das analoge Signale in einen digitalen Datenstrom, der dann weiterverarbeitet oder gespeichert werden kann}
}

\newglossaryentry{notch}{
    name=Notch-Filter,
    first={\textit{Notch-Filter}},
    description={Filter, der gezielt eine bestimmte Frequenz (z.B. 50 Hz Netzrauschen) unterdrückt}
}

\newglossaryentry{icalang}{
    name=Independent Component Analysis,
    first={\textit{Independent Component Analysis}},
    description={Independent Component Analysis ist ein mathematisches Verfahren zur Trennung überlagerter Signalquellen im EEG}
}

\newglossaryentry{bandpass}{
    name=Bandpassfilter,
    first={\textit{Bandpassfilter}},
    description={Filter, der nur Signale innerhalb eines bestimmten Frequenzbereichs durchlässt}
}

\newglossaryentry{cutoff}{
    name=Cut-off-Frequenz,
    first={\textit{Cut-off-Frequenz}},
    description={Grenzfrequenz, ab der ein Filter Signale durchlässt oder unterdrückt}
}

\newglossaryentry{p300}{
    name=P300,
    first={\textit{P300}},
    description={ERP-Komponente, die etwa 300 ms nach einem seltenen oder bedeutsamen Stimulus auftritt}
}

\newglossaryentry{elektrookulogramm}{
    name=Elektrookulogramm,
    first={\textit{Elektrookulogramm}},
    description={Elektrookulogramme sind Messungen von Augenbewegungen zur Identifikation von Artefakten im EEG}
}

\newglossaryentry{eventmarker}{
    name=Ereignismarker,
    first={\textit{Ereignismarker}},
    description={Zeitliche Markierung im EEG-Signal zur Kennzeichnung eines Stimulus oder Ereignisses}
}

\newglossaryentry{oddball}{
    name=Oddball-Aufgabe,
    first={\textit{Oddball-Aufgabe}},
    description={Experimentelles Paradigma, bei dem seltene Zielreize zwischen häufigen Standardreizen präsentiert werden, um ERPs zu untersuchen}
}

\newglossaryentry{epikrImpulse}{
    name=Epikriptische Impulse,
    first={\textit{Epikriptische Impulse}},
    description={Feine Sinnesreize wie Berührung oder Druck, die präzise wahrgenommen und genau lokalisiert werden können}
}

\newglossaryentry{ProtoImpulse}{
    name=Protopathische Impulse,
    first={\textit{Protopathische Impulse}},
    description={Grobe Sinnesreize wie Schmerz oder Temperatur, die weniger genau unterschieden werden können}
}

\newglossaryentry{wernick}{
    name=Wernicke-Sprachzentrum,
    first={\textit{Wernicke-Sprachzentrum}},
    description={Gehirnbereich, der für das Verstehen von gesprochener und geschriebener Sprache zuständig ist}
}

\newglossaryentry{oligo}{
    name=Oligodendrozyten,
    first={\textit{Oligodendrozyten}},
    description={Spezialisierte Zellen im Nervensystem, die Nervenfasern mit einer schützenden Myelinschicht umhüllen}
}


\newglossaryentry{rückenmark}
    {
        name=Rückenmark,
        first={\textit{Rückenmark}},
        description={Das Rückenmark ist ein Teil des zentralen Nervensystems, der sich im Wirbelkanal befindet und als Hauptleitung für Nervenimpulse zwischen Gehirn und Körper dient}
    }

\newglossaryentry{peripheres Nervensystem}
    {
        name=peripheres Nervensystem,
        first={\textit{peripheres Nervensystem}},
        description={Das periphere Nervensystem umfasst alle Nerven außerhalb des zentralen Nervensystems (Gehirn und Rückenmark) und verbindet diese mit den Organen, Muskeln und Haut}
    }
    

\newglossaryentry{spährischer Spline-Interpolation}
    {
    name=spährische Spline-Interpolation,
    first={\textit{spährischer Spline-Interpolation}},
    description={Die sphärische Spline-Interpolation im EEG ist ein mathematisches Verfahren zur raumbezogenen Schätzung und Glättung von elektrischen Potenzialen auf einer Kugeloberfläche.}
    }
 
    \newglossaryentry{no-switzerland-principle}
    {
    name=grand-average,
    first={\textit{no-switzerland-principle}},
    description={123456}
    }
 
   \newglossaryentry{grand-average}
    {
    name=grand-average,
    first={\textit{grand-average}},
    description={123456}
    }

    \newglossaryentry{Power Specral Density}
    {
    name=Power Specral Density,
    first={\textit{Power Specral Density}},
    description={Die spektrale Leistungsdichte (PSD) im EEG beschreibt, wie sich die Signalenergie oder -leistung eines Elektroenzephalogramms auf verschiedene Frequenzen verteilt. Sie wandelt rohe, komplexe Zeitreihendaten in den Frequenzbereich um und macht so Grundrhythmen wie Delta-, Theta-, Alpha- und Betawellen sichtbar.}
    }

    \newglossaryentry{Welch-Methode}
    {
    name=Welch-Methode,
    first={\textit{Welch-Methode}},
    description={Die Welch-Methode ist ein Verfahren zur Bestimmung der spektralen Leistungsdichte (PSD) eines Signals. Dabei wird das EEG-Signal in überlappende Abschnitte geteilt, deren Spektren berechnet und anschließend gemittelt werden. Dadurch entsteht eine stabilere Darstellung der Frequenzanteile des EEG-Signals.}
    }

    \newglossaryentry{openbci-board}
    {
    name=OpenBCI-Board,
    first={\textit{OpenBCI-Board}},
    description={}
    }

    \newglossaryentry{brainflow}{
        name = BrainFlow,
        first={\textit{BrainFlow}},
        description={}
        }
    \newglossaryentry{abtastrate}{
        name = Abtastrate,
        first={\textit{Abtastrate}},
        description={}
        }
    

\newglossaryentry{mne}{
    name={MNE},
    first={\textit{MNE}},
    description={}
}

    \newglossaryentry{stimulus}{
    name={Stimulus},
    first={\textit{Stimulus}},
    description={Ein Stimulus (Mehrzahl Stimuli) ist ein Reiz oder Anreiz, der eine bestimmte körperliche, psychische oder wirtschaftliche Reaktion auslöst.}
}
    

%hier kommen die Einträge für das Abkürzungsverzeichnis rein
%so:
%\newacronym{label}{abkürzung}{lang}

\newacronym{csv}{CSV}{Comma-Separated Values, \glqq durch Komma getrennte Werte\grqq}
\newacronym{bdf}{BDF}{BioSemi Data Format \glqq BioSemi-Datenformat\grqq}

\newacronym{eeg}{EEG}{Elektroenzephalographie oder Elektroenzephalogramm, Mehrzahl EEGs}

\newacronym{bci}{BCI}{Brain-Computer-Interface, \glqq Gehirn-Computer-Interface\grqq, Mehrzahl BCIs}

\newacronym{epsp}{EPSP}{exzitatorisches / erregendes postsynaptisches Potenzial, Mehrzahl EPSPs}

\newacronym{ipsp}{IPSP}{inhibitorisches / hemmendes postsynaptisches Potenzial, Mehrzahl IPSPs}

\newacronym{adc}{ADC}{\gls{analog-digital-converter}, \glqq Analog-Digital-Umsetzer\grqq}

\newacronym{erps}{ERPs}{Mehrzahl von \acrshort{erp}}

\newacronym{erp}{ERP}{Event-related potentials, \glqq ereignisbezogene Potenziale\grqq, Mehrzahl ERPs}

\newacronym{dc}{DC}{Direct Current, \glqq Gleichstrom\grqq}

\newacronym{eog}{EOG}{Elektrookulographie}

\newacronym{cms}{CMS}{Common-Sense, \glqq gesunder Menschenverstand \grqq, hier ist eine CMS-Elektrode eine Referenzelektrode\grqq}

\newacronym{drl}{DRL}{Driven-Right-Leg, \glqq Rechte-Bein-Ansteuerung\grqq}

\newacronym{ica}{ICA}{Independent Component Analysis, \glqq Unabhängigkeitsanalyse\grqq}

\newacronym{ekp}{EKPs}{ereigniskorrelierte Potenziale}


%gkleiche abstände zwischen den Absätzen
\newglossarystyle{acrostyle}{%
  \setglossarystyle{long}%
  \setlength{\glsdescwidth}{0.7\textwidth}%
}



%Formatierung und Import Bibliography

\usepackage[affil-it]{authblk}	
%Zitierweise nach CZG: Kurzbeleg in der Fußnote "Vgl. Familienname (Jahr): S. x."
%Quellenverzeichnis: alle Autor:innen (maxbibnames=99); im Kurzbeleg bei 3-5 Autor:innen
%beim ersten Zitat alle Namen, danach "Erstautor et al." (siehe \AtEveryCitekey unten);
%bei mehr als 5 Autor:innen von Anfang an "et al.".
\usepackage[datamodel=datamodel,backend=biber,style=authoryear,autocite=footnote,
  maxbibnames=99,minbibnames=99,maxcitenames=5,mincitenames=1,
  uniquename=false,uniquelist=false,citetracker=true,
  urldate=short,lasteditdate=long,dateabbrev=false]{biblatex}
\addbibresource{literatur.bib}
\setlength{\bibitemsep}{1.5em}

% slash zwischen namen 
\DeclareDelimFormat{multinamedelim}{\space\slash\space}
\DeclareDelimFormat{finalnamedelim}{\space\slash\space}
%Vor und nach name richtige reinfolge
\DeclareNameAlias{sortname}{family-given}   
\DeclareNameAlias{default}{family-given}    

\DefineBibliographyStrings{ngerman}{
  urlseen = {Abrufdatum},
  andothers = {et\addabbrvspace al\adddot}
}

%URL ohne vorangestelltes "url:"; Abrufdatum als [Abrufdatum TT.MM.JJJJ]
\DeclareFieldFormat{url}{\url{#1}}
\DeclareFieldFormat{pages}{S\adddot\addnbspace#1}
\DeclareFieldFormat{edition}{\ifinteger{#1}{#1\adddot\addspace Aufl\adddot}{#1}}

%Kurzbeleg (\cite): Druckquellen "Familienname (Jahr)", Internetquellen
%"Autor/Institution (Abrufdatum)", Lexika "Lexikon: Stichwort (Abrufdatum)".
%Gibt es von derselben Institution mehrere Seiten mit gleichem Abrufdatum, steht
%zusätzlich das Feld shorttitle hinter dem Namen.
\newbibmacro{czg:abrufdatum}{%
  \printtext[parens]{\mkbibdateshort{urlyear}{urlmonth}{urlday}}}
\newbibmacro{czg:cite}{%
  \ifentrytype{onlinelexikon}
    {\printtext{\thefield{journaltitle}\addcolon\space\thefield{title}}%
     \setunit{\addspace}\usebibmacro{czg:abrufdatum}}
    {\printnames{labelname}%
     \ifboolexpr{test {\ifentrytype{onlinedokument}} and not test {\iffieldundef{urlyear}}}
       {\iffieldundef{shorttitle}{}{\printtext{\addcolon\space\thefield{shorttitle}}}%
        \setunit{\addspace}\usebibmacro{czg:abrufdatum}}
       {\setunit{\addspace}\printtext[parens]{\printlabeldateextra}}}}

\DeclareCiteCommand{\cite}
  {\usebibmacro{prenote}}
  {\usebibmacro{citeindex}\usebibmacro{czg:cite}}
  {\multicitedelim}
  {\usebibmacro{postnote}}
\DeclareDelimFormat{postnotedelim}{\addcolon\space}
\DeclareFieldFormat{postnote}{#1}
\renewcommand*{\multicitedelim}{\addsemicolon\space}

%3-5 Autor:innen: beim ersten Zitat alle Namen, danach "et al."
%Endet die Autorenliste im Literatureintrag mit "and others" (mehr Autor:innen als
%eingetragen), steht von Anfang an "Erstautor et al."
\AtEveryCitekey{\ifciteseen{\defcounter{maxnames}{2}}{\defcounter{maxnames}{5}}%
  \ifandothers{labelname}{\defcounter{maxnames}{1}}{}}

%Datum ohne Leerzeichen: TT.MM.JJJJ
\DefineBibliographyExtras{ngerman}{%
  \protected\def\mkbibdateshort#1#2#3{%
    \iffieldundef{#3}{}{\mkdayzeros{\thefield{#3}}.}%
    \iffieldundef{#2}{}{\mkmonthzeros{\thefield{#2}}.}%
    \iffieldundef{#1}{}{\mkyearzeros{\thefield{#1}}}}}

%\DeclareFieldFormat{note}{\mkbibbrackets{\bibstring{urlseen}\addcolon\space#1}}

\DeclareFieldFormat{urldate}{\mkbibbrackets{\bibstring{urlseen}\space#1}}

\DeclareFieldFormat{lasteditdate}{\mkbibparens{#1}}

\DeclareFieldFormat{year}{\mkbibparens{#1}}

\DeclareFieldFormat{volume}{#1}






%citation commands
\DeclareCiteCommand{\vgl}[\mkbibfootnote]{%
	\usebibmacro{citeindex}%
	\iffieldundef{prenote}
		{\printtext{Vgl\adddot\space}}
		{\printfield{prenote}\setunit{\prenotedelim}}%
}
{%
	% Ausgabe abhängig vom Eintragstyp
	\ifentrytype{onlinedokument}
		{\printnames{labelname}\addspace\bibhyperref{\printurldate}}%
		{%
					\ifentrytype{onlinelexikon}
						{\printfield{journaltitle}\adddot\addspace\bibhyperref{\printfield{title}}}%
						{%
							\ifentrytype{zeitschriftartikel}
								{\printnames[family]{labelname}\addspace\bibhyperref{\printfield{year}}}%
								{\printnames[family]{labelname}\addspace\bibhyperref{\printfield{year}}}%
						}%
		}%
}
{\multicitedelim}
{%
	\iffieldundef{postnote}
		{\adddot}
		{\addcolon\addspace\printfield{postnote}\adddot}%
	\nopunct%
}

\DeclareCiteCommand{\vglHier}[]{%
	\usebibmacro{citeindex}%
	\iffieldundef{prenote}
		{\printtext{Vgl\adddot\space}}
		{\printfield{prenote}\setunit{\prenotedelim}}%
}
{%
	% Ausgabe abhängig vom Eintragstyp
	\ifentrytype{onlinedokument}
		{\printnames{labelname}\addspace\bibhyperref{\printurldate}}%
		{%
					\ifentrytype{onlinelexikon}
						{\printfield{journaltitle}\adddot\addspace\bibhyperref{\printfield{title}}}%
						{%
							\ifentrytype{zeitschriftartikel}
								{\printnames[family]{labelname}\addspace\bibhyperref{\printfield{year}}}%
								{\printnames[family]{labelname}\addspace\bibhyperref{\printfield{year}}}%
						}%
		}%
}
{\multicitedelim}
{%
	\iffieldundef{postnote}
		{\adddot}
		{\addcolon\addspace\printfield{postnote}\adddot}%
	\nopunct%
}
	


\DeclareCiteCommand{\ebd}[\mkbibfootnote]
{\usebibmacro{citeindex}%
\iffieldundef{prenote}
{\printtext{Vgl\adddot}}
{\printfield{prenote}\setunit{\prenotedelim}}%
}
{
	\printtext{ebd\adddot}
	\iffieldundef{postnote}
	{}
	{\addspace\printfield{postnote}\adddot}
}
{}
{}




%Einbinden von Quelltext in Python (https://tex.stackexchange.com/questions/83882/how-to-highlight-python-syntax-in-latex-listings-lstinputlistings-command)
% Default fixed font does not support bold face
\DeclareFixedFont{\ttb}{T1}{txtt}{bx}{n}{12} % for bold
\DeclareFixedFont{\ttm}{T1}{txtt}{m}{n}{12}  % for normal

% Custom colors
\usepackage{color}
\definecolor{deepblue}{rgb}{0,0,0.5}
\definecolor{deepred}{rgb}{0.6,0,0}
\definecolor{deepgreen}{rgb}{0,0.5,0}

\usepackage{listings}

% Python style for highlighting
\newcommand\pythonstyle{\lstset{
language=Python,
basicstyle=\ttm,
otherkeywords={self},             % Add keywords here
keywordstyle=\ttb\color{deepblue},
emph={MyClass,__init__},          % Custom highlighting
emphstyle=\ttb\color{deepred},    % Custom highlighting style
stringstyle=\color{deepgreen},
frame=tb,                         % Any extra options here
showstringspaces=false            % 
}}


% Python environment
\lstnewenvironment{python}[1][]
{
\pythonstyle
\lstset{#1}
}
{}

% Python for external files
\newcommand\pythonexternal[2][]{{
\pythonstyle
\lstinputlisting[#1]{#2}}}

% Python for inline
\newcommand\pythoninline[1]{{\pythonstyle\lstinline!#1!}}

\usepackage{eurosym}
